\section{Структура данных и пайплайн подготовки датасета}
\label{sec:data-pipeline}

\subsection{Форматы входных данных}
В проекте используются три формата (табл.~\ref{tab:formats}).

\begin{table}[h]
\centering
\caption{Форматы данных и их назначение}
\label{tab:formats}
\begin{tabular}{|p{2.5cm}|p{3cm}|p{8cm}|}
\hline
\textbf{Формат} & \textbf{Расширение} & \textbf{Назначение} \\
\hline
STEP & .stp, .step & Точное B-Rep представление, используется для автоматической разметки и для детерминированных вычислений в CAD \\
\hline
STL & .stl & Триангулированная сетка, вход нейросети (MeshCNN) \\
\hline
JSON & .json & Структурированное описание истории построения (тип операции, параметры, ссылки на грани) \\
\hline
\end{tabular}
\end{table}

\subsection{Формирование и очистка обучающей выборки}
Исходный массив содержал 82~418 моделей, распределённых по папкам. Для каждой модели требовалось наличие как минимум двух последовательных шагов (начальная и конечная геометрия) и JSON-файла с описанием операции перехода. Процесс очистки включал следующие шаги:

\begin{enumerate}
    \item Удаление папок, в которых отсутствует JSON-файл (262 модели).
    \item Удаление моделей, у которых нет ни STL, ни STEP (7~417 моделей).
    \item Отбраковка моделей с единственной парой STL+STEP (нет соседних шагов) - 15~020 моделей.
    \item Ручная проверка аномалий: два JSON в одной папке (18 случаев), пустые STL, нечитаемые STEP.
\end{enumerate}

В результате получен рабочий набор из \textbf{21~464 моделей}, каждая из которых содержит полную историю построения и может быть использована для обучения. Статистика приведена в таблице~\ref{tab:data-stats}.

\begin{table}[h]
\centering
\caption{Статистика подготовки данных}
\label{tab:data-stats}
\begin{tabular}{|l|r|p{6cm}|}
\hline
\textbf{Показатель} & \textbf{Значение} & \textbf{Смысл для проекта} \\
\hline
Подготовлено моделей & 82 418 & Сырой массив, отражающий потенциал масштабирования \\
\hline
Прошли первичную обработку & 44 181 & Модели, достигшие стадии первичной машинной обработки \\
\hline
Два JSON & 18 & Структурная неоднозначность метаданных \\
\hline
Нет JSON & 262 & Невозможность supervised-разметки по операциям \\
\hline
Нет STL/STEP & 7 417 & Потеря геометрической основы для обучения \\
\hline
Только одна пара STL+STEP & 15 020 & Непригодность для обучения переходам и шагам \\
\hline
\textbf{Итоговый рабочий набор} & \textbf{21 464} & Модели с полным составом данных \\
\hline
Подвыборка для fillet & 1 726 & Фактический обучающий массив для частного сегментационного эксперимента \\
\hline
\end{tabular}
\end{table}

\subsection{Нормализация и препроцессинг}
Для каждого STL-файла, попадающего в обучение или инференс, выполняются следующие операции:
\begin{itemize}
    \item \textbf{Центрирование}: все вершины сдвигаются так, чтобы центроид модели оказался в начале координат.
    \item \textbf{Масштабирование}: модель масштабируется к единичному размеру (максимальная координата становится равной 1).
    \item \textbf{Приведение числа рёбер}: с помощью процедур subdivision/decimation количество рёбер фиксируется на значении 2250 (требование MeshCNN).
    \item \textbf{Удаление вырожденных граней}: треугольники с нулевой площадью или слишком острыми углами исключаются.
\end{itemize}

Параметры операций (глубина, радиус, угол) также нормализуются: для каждого параметра вычисляется среднее и стандартное отклонение по обучающей выборке, после чего применяется z-нормировка.

\subsection{Автоматическая разметка датасета}
Разметка рёбер для задачи сегментации выполняется автоматически на основе JSON-описания и STEP-файлов соседних шагов. Алгоритм:
\begin{enumerate}
    \item Для пары соседних шагов (до операции и после операции) загружаются STEP-модели.
    \item Вычисляются грани, которые изменились (добавлены, удалены, модифицированы).
    \item По JSON определяется тип операции, породившей изменение.
    \item Для каждого треугольника в финальной STL-сетке определяется, принадлежит ли он изменённой грани; если да, то все рёбра этого треугольника получают метку, соответствующую типу операции.
    \item Результат сохраняется в формате Colored OBJ, где рёбра кодируются тройками (индекс вершины 1, индекс вершины 2, метка класса).
\end{enumerate}
Такой подход позволяет получить размеченный датасет без ручной аннотации, однако он чувствителен к качеству сравнения STEP-моделей и может вносить ошибки при сложных топологических изменениях.

